\documentclass{article}
\usepackage{graphicx} % Required for inserting images
\usepackage{amsmath}  % Required for align environment
\usepackage[margin=1in]{geometry}


\title{ELCT 432 Assignment \#1}
\author{Jacob C. Vaught}
\date{August 24th 2023}

\begin{document}

\maketitle

\section*{Problem \#1}
\subsection*{Problem Statement}
Show that the sum of two discrete periodic signals is periodic, whereas the sum of two continuous periodic signals is not necessarily periodic. Under what condition is the sum of two continuous periodic signals periodic?

\subsection*{Part 1: Sum of Two Discrete Periodic Signals}

Let \( x(t) \) and \( y(t) \) be two discrete periodic signals with periods \( T_1 \) and \( T_2 \), respectively. Their sum \( z(t) = x(t) + y(t) \) is also periodic with period \( T \) equal to the least common multiple (LCM) of \( T_1 \) and \( T_2 \):

\begin{align*}
z(t + T) & = x(t + \text{LCM}(T_1, T_2)) + y(t + \text{LCM}(T_1, T_2)) \\
 & = x(t) + y(t) \\
 & = z(t)
\end{align*}

\subsection*{Part 2: Sum of Two Continuous Periodic Signals}

For continuous periodic signals \( x(t) \) and \( y(t) \) with periods \( T_1 \) and \( T_2 \), respectively, we assume that the ratio of their periods is rational: \( \frac{T_1}{T_2} = \frac{m}{n} \), where \( m \) and \( n \) are integers with no common divisors.
\newline \newline We then find integers \( k_1 \) and \( k_2 \) such that:
\begin{align*}
k_1 T_1 = k_2 T_2
\end{align*}
The period \( T \) of the sum of the signals will be:
\begin{align*}
T = \text{LCM}(T_1, T_2) = k_1 T_1 = k_2 T_2
\end{align*}
So the sum of the two signals is:
\begin{align*}
z(t + T) & = x(t + k_1 T_1) + y(t + k_2 T_2) \\
 & = x(t) + y(t) \\
 & = z(t)
\end{align*}
Therefore, the sum of two continuous periodic signals is periodic if the ratio of their periods is rational.


\newpage
\section*{Problem \#2}
\subsection*{Problem Statement}
The objective of this analysis is to ascertain the periodicity of the given signals. Specifically, we aim to:
\begin{itemize}
    \item Determine whether each signal is periodic.
    \item For signals that are found to be periodic, identify their period.
\end{itemize}
The signals under consideration are as follows:
\begin{enumerate}
    \item \( x(t) = \sin(4000\pi t) + \cos(1000\pi t) \)
    \item \( x(t) = \sin(4000\pi t) + \cos(11000t) \)
    \item \( x[n] = \sin(4000\pi n) + \cos(11000\pi n) \)
    \item \( x[n] = \sin(4000\pi n) + \cos(11000n) \)
\end{enumerate}

\subsection*{Signal 1: \( x(t) = \sin(4000\pi t) + \cos(1000\pi t) \)}
Both \( \sin(4000\pi t) \) and \( \cos(1000\pi t) \) are periodic with periods \( T_1 = \frac{1}{2000} \) and \( T_2 = \frac{1}{500} \), respectively. 
The sum is also periodic with period \( T = \text{LCM}(T_1, T_2) = \text{LCM}\left(\frac{1}{2000}, \frac{1}{500}\right) = \frac{1}{500} \).

\subsection*{Signal 2: \( x(t) = \sin(4000\pi t) + \cos(11000t) \)}
The function \( \sin(4000\pi t) \) is periodic with \( T_1 = \frac{1}{2000} \). However, \( \cos(11000t) \) is not periodic since the frequency is not a multiple of \( \pi \). Therefore, the sum is non-periodic.

\subsection*{Signal 3: \( x[n] = \sin(4000\pi n) + \cos(11000\pi n) \)}
Both \( \sin(4000\pi n) \) and \( \cos(11000\pi n) \) are discrete periodic signals with fundamental periods \( N_1 = 1 \) and \( N_2 = 1 \), respectively. Their sum is also periodic with \( N = \text{LCM}(N_1, N_2) = \text{LCM}(1, 1) = 1 \).

\subsection*{Signal 4: \( x[n] = \sin(4000\pi n) + \cos(11000n) \)}
The function \( \sin(4000\pi n) \) is periodic with \( N_1 = 1 \). However, \( \cos(11000n) \) is not periodic in the general sense of discrete-time signals. Thus, the sum is non-periodic.

\newpage
\section*{Problem \#3}

\subsection*{Problem Statement}
Prove that a system \( \mathcal{F} \) is linear if and only if it satisfies the following conditions:
\begin{enumerate}
    \item Homogeneity: For all input signals \( x(t) \) and all real numbers \( \alpha \), \( \mathcal{F}\{\alpha x(t)\} = \alpha \mathcal{F}\{x(t)\} \).
    \item Additivity: For all input signals \( x_1(t) \) and \( x_2(t) \), \( \mathcal{F}\{x_1(t) + x_2(t)\} = \mathcal{F}\{x_1(t)\} + \mathcal{F}\{x_2(t)\} \).
\end{enumerate}

\subsection*{Solution}

\subsubsection*{Step 1: If \( \mathcal{F} \) is Linear, Then It is Homogeneous and Additive}

Assume that \( \mathcal{F} \) is a linear system. Then for any \( x_1(t) \) and \( x_2(t) \), and any scalars \( \alpha \) and \( \beta \):

\[
\mathcal{F}\{\alpha x_1(t) + \beta x_2(t)\} = \alpha \mathcal{F}\{x_1(t)\} + \beta \mathcal{F}\{x_2(t)\}
\]
Taking \( \beta = 0 \), 
\[
\mathcal{F}\{\alpha x_1(t)\} = \alpha \mathcal{F}\{x_1(t)\}
\]
which proves homogeneity.
\newline
\newline Similarly, taking \( \alpha = \beta = 1 \):
\[
\mathcal{F}\{x_1(t) + x_2(t)\} = \mathcal{F}\{x_1(t)\} + \mathcal{F}\{x_2(t)\}
\]
which proves additivity.

\subsubsection*{Step 2: If \( \mathcal{F} \) is Homogeneous and Additive, Then It is Linear}

Assume that \( \mathcal{F} \) is homogeneous and additive. Then for any \( x_1(t) \) and \( x_2(t) \), and any scalars \( \alpha \) and \( \beta \):

\[
\mathcal{F}\{\alpha x_1(t)\} = \alpha \mathcal{F}\{x_1(t)\} \quad \text{(by homogeneity)}
\]
\[
\mathcal{F}\{\beta x_2(t)\} = \beta \mathcal{F}\{x_2(t)\} \quad \text{(by homogeneity)}
\]
Since \( \mathcal{F} \) is also additive:
\[
\mathcal{F}\{\alpha x_1(t) + \beta x_2(t)\} = \mathcal{F}\{\alpha x_1(t)\} + \mathcal{F}\{\beta x_2(t)\}
\]
Combining these:
\[
\mathcal{F}\{\alpha x_1(t) + \beta x_2(t)\} = \alpha \mathcal{F}\{x_1(t)\} + \beta \mathcal{F}\{x_2(t)\}
\]
Thus, \( \mathcal{F} \) is linear.
\newline \newline Therefore, it is proven that \( \mathcal{F} \) is linear if and only if it is homogeneous and additive.


\newpage
\section*{Problem \#4}
\subsection*{Problem Statement}
\textbf{Objective:} To analyze the properties of a specific amplitude modulation system, including its linearity and time-invariance.
\newline \newline
\textbf{System Description:} 
The system under consideration is defined by the input-output relation:
\[
y(t) = x(t) \cos(2\pi f_{\text{carrier}} t)
\]
where \( f_{\text{carrier}} \) is a constant representing the carrier frequency. This system is commonly referred to as a modulator.
\newline \newline
\textbf{Questions:}

\begin{enumerate}
    \item Is this system linear?
    \item Is this system time-invariant?
\end{enumerate}
\newline 
\textbf{Instructions:} Provide a rigorous analysis to answer each question, including any necessary mathematical derivations.

\subsection*{Solution}
\textbf{Linearity}
\newline
A system is considered linear if it satisfies the principles of superposition and homogeneity. Mathematically, if \( y_1(t) \) and \( y_2(t) \) are the outputs corresponding to inputs \( x_1(t) \) and \( x_2(t) \) respectively, then for any constants \( a \) and \( b \), the system must satisfy:

\begin{align*}
a y_1(t) + b y_2(t) &= a x_1(t) \cos(2\pi f_{\text{carrier}} t) + b x_2(t) \cos(2\pi f_{\text{carrier}} t) \nonumber \\
&= (a x_1(t) + b x_2(t)) \cos(2\pi f_{\text{carrier}} t)
\end{align*}
\newline
Notice that the last term on the right-hand side is exactly the output for an input \( a x_1(t) + b x_2(t) \). Therefore, the system is linear.
\newline \newline
\textbf{Time Invariance}
\newline
A system is time-invariant if shifting the input signal in time produces an identical shift in the output signal. Mathematically, if \( y(t) \) is the output corresponding to an input \( x(t) \), then for a time shift \( \tau \):
\newline
\begin{align*}
y(t - \tau) &= x(t - \tau) \cos(2\pi f_{\text{carrier}} (t - \tau)) \\
&= x(t - \tau) \cos(2\pi f_{\text{carrier}} t - 2\pi f_{\text{carrier}} \tau)
\end{align*}
\newline
The term \( 2\pi f_{\text{carrier}} \tau \) in the cosine argument indicates that the system is not time-invariant, as an additional phase term is introduced when the input is shifted.

\subsection*{Conclusion}
\begin{itemize}
    \item The system is linear.
    \item The system is not time-invariant.
\end{itemize}

\newpage
\section*{Problem \#5}
\subsection*{Problem Statement}
Find the Euler-form representation for the following signals:

\begin{itemize}
    \item \( x(t) = \cos(2\pi t) + \cos(4\pi t) \)
    \item \( x(t) = \cos(2\pi t) - \cos(4\pi t + \frac{\pi}{3}) \)
    \item \( x(t) = 2\cos(2\pi t) - \sin(4\pi t) \)
\end{itemize}

\subsection*{Solution}
\textbf{First Signal: \( x(t) = \cos(2\pi t) + \cos(4\pi t) \)}
\newline
Using Euler's formula, rewrite \( x(t) \) as:
\newline
\[
x(t) = \frac{e^{j 2\pi t} + e^{-j 2\pi t}}{2} + \frac{e^{j 4\pi t} + e^{-j 4\pi t}}{2}
\]
\newline \newline
\textbf{Second Signal: \( x(t) = \cos(2\pi t) - \cos(4\pi t + \frac{\pi}{3}) \)}
\newline
Again, using Euler's formula, this signal becomes:
\newline
\[
x(t) = \frac{e^{j 2\pi t} + e^{-j 2\pi t}}{2} - \frac{e^{j (4\pi t + \frac{\pi}{3})} + e^{-j (4\pi t + \frac{\pi}{3})}}{2}
\]
\newline \newline
\textbf{Third Signal: \( x(t) = 2\cos(2\pi t) - \sin(4\pi t) \)}
\newline
For this signal, the Euler-form representation is:
\newline
\[
x(t) = 2 \left( \frac{e^{j 2\pi t} + e^{-j 2\pi t}}{2} \right) - \frac{e^{j 4\pi t} - e^{-j 4\pi t}}{2j}
\]
\end{document}